\documentclass[aps, superscriptaddress, floatfix, onecolumn]{revtex4-2}
\usepackage{amsmath}
\usepackage{amssymb}
\usepackage{amsthm}
\usepackage{mathtools}
\usepackage{graphicx}
\usepackage{braket}

% For a list of named colors, see https://www.overleaf.com/learn/latex/Using_colors_in_LaTeX
\usepackage{color}
\usepackage[dvipsnames]{xcolor}

\newcommand{\result}[1]{\textcolor{blue}{#1}}



\newcommand{\CCQ}{Center for Computational Quantum Physics, Flatiron Institute, 162 5th Avenue, New York, NY 10010, USA}

\newtheorem{remark}{Remark}
\newcommand{\ms}[1]{\textcolor{blue}{[MS: #1]}}
\newtheorem{theorem}{Theorem}

\begin{document}

\title{Gaussian Fermions}

\author{E.\ Miles Stoudenmire}
\affiliation{\CCQ}

\maketitle

\section{Fermionic Vector Spaces}

Fermions are conventionally defined in terms of creation and annihilation operators. Here we will instead start from fermionic vector spaces or state spaces and the properties which make them ``fermionic.'' In the next section, we will see that the more familiar properties of creation and annihilation operators follow from the properties of the states.

The simplest fermionic vector space $V$ has the basis 
\begin{align}
    \mathcal{B}(V) = \{ \ket{0}, \ket{1} \}
\end{align}
which is the basis for a single ``spinless fermion,'' with $\ket{0}$ being the vacuum or empty state and $\ket{1}$ being the filled state or state of one fermion. So far the fermionic nature is not very obvious except  that we can only have at most one fermion. 

Many-body spaces can be formed by taking Cartesian products of spaces (tensor products of basis vectors) so we could consider a space $V_1$ with basis $\{ \ket{0}_1, \ket{1}_1 \}$ and space $V_2$ with basis $\{ \ket{0}_2, \ket{1}_2 \}$, then the space $V_1 \times V_2$ has a basis
\begin{align}
    \mathcal{B}(V_1 \times V_2) = \{ \ket{0}_1 \otimes \ket{0}_2,\  \ket{1}_1 \otimes \ket{0}_2, \  \ket{0}_1 \otimes \ket{1}_2,\  \ket{1}_1  \otimes \ket{1}_2 \} \ .
\end{align}
A generic vector (or ``state'') in this space has the form $\ket{\psi} = a \ket{0}_1 \otimes \ket{0}_2 + b \ket{0}_1 \otimes \ket{1}_2 + c \ket{1}_1 \otimes \ket{0}_2 + d \ket{1}_1 \otimes \ket{1}_2 $.

The defining property of fermions is how one identifies vectors in the space $V_1 \times V_2$  with vectors in the space $V_2 \times V_1$. We can define this identification, or vector space isomorphism, by defining it on the basis vectors. It is given by:
\begin{align}
\ket{0}_1 \otimes \ket{0}_2 & \leftrightarrow {\color{ForestGreen} +}\ket{0}_2 \otimes \ket{0}_1 \nonumber \\
\ket{1}_1 \otimes \ket{0}_2 & \leftrightarrow {\color{ForestGreen} +}\ket{0}_2 \otimes \ket{1}_1 \nonumber \\
\ket{0}_1 \otimes \ket{1}_2 & \leftrightarrow {\color{ForestGreen} +}\ket{1}_2 \otimes \ket{0}_1  \nonumber \\
\ket{1}_1 \otimes \ket{1}_2 & \leftrightarrow {\color{BrickRed} -}\ket{1}_2 \otimes \ket{1}_1  \ \ .
\end{align}
Note the minus sign in the last line. Plus signs are shown for emphasis.

We can summarize the vector space isomorphism as
\begin{align}
     \ket{n_1}_1 \otimes \ket{n_2}_2 \ \leftrightarrow\ (-1)^{n_1 \cdot n_2} \ket{n_2}_2 \otimes \ket{n_1}_1 
\end{align}
where $n_1, n_2 \in \{0,1\}$.
The factor in front of this isomorphism rule means the \emph{ordering} of fermionic basis vectors matters: if two odd basis vectors are permuted, the basis vector changes by a minus sign. 
\footnote{In some mathematical literature, an isomorphism of this type is described as ``having a $\mathbb{Z}_2$ grading'', meaning that one assigns non-trivial properties based on the parity or grading of the basis, or having a ``braided'' or ``graded'' vector space isomorphism.}

Note that tensor products of basis states such as $\ket{1}_i \otimes \ket{1}_j$ are commonly written in the more compact notation $\ket{1}_i \ket{1}_j$ or even $\ket{1_i 1_j}$. Tensor product symbols ``$\otimes$'' are generally dropped when understood from context, and will be dropped frequently in the rest of these notes.


\section{Fermionic Operators}

For the simplest case of a ``spinless fermion'' space $V$ with basis $\mathcal{B}(V) = \{ \ket{0}, \ket{1} \}$, we can define the following operators:
\begin{align}
\hat{n}\ &= \ \ket{1} \bra{1} \\ 
\hat{c}\ &= \ \ket{0} \bra{1} \\ 
\hat{c}^\dagger &= \ \ket{1} \bra{0}  \ \ .
\end{align}
Together with the identity operator, this set defines a complete operator basis.

The $\hat{n}$ operator is known as the ``number operator''
since it ``detects'' the filling of a state:
\begin{align}
\hat{n} \ket{0}  &= 0 \cdot \ket{0}  \nonumber \\ 
\hat{n} \ket{1}  &= 1 \cdot \ket{1}  \ \ .
\end{align}

The $\hat{c}$ and $\hat{c}^\dagger$ are known as the ``annihilation'' and ``creation'' operators, respectively. This is simply because they map a filled state to an empty state or vice versa:
\begin{align}
\hat{c} \ket{1}  &= \ket{0}  \\ 
\hat{c}^\dagger \ket{0}  &= \ket{1}  \ \ .
\end{align}

One can of course define tensor products of operators in the usual way, such as $\hat{c}_1 \hat{c_2}$ or $\hat{c}^\dagger_2  \hat{c}_1$.
\footnote{When an operator like $\hat{c}_1$ is written in isolation, it actually stands for $\hat{c}_1 \otimes \hat{I}_2$. Products of operators such as $\hat{c}^\dagger_1 \hat{c}_2$ are actually shorthand for $(\hat{c}^\dagger_1 \otimes \hat{I}_2) (\hat{I}_1 \otimes \hat{c}_2) = \hat{c}^\dagger_1 \otimes \hat{c}_2$.
If the operators act on the same space, they ``stack'', meaning for example that $(\hat{c}^\dagger_1 \otimes \hat{I}_2) (\hat{c}_1 \otimes \hat{I}_2) = (\hat{c}^\dagger_1 \hat{c}_1) \otimes \hat{I}_2$. All of this is easier to see with tensor diagrams!}
These operators act on the spaces $V_1 \times V_2$ or $V_2 \times V_1$, respectively.

Now comes a crucial point: the creation and annihilation operators have non-trivial commutation relations, and these are entirely a consequence of the state properties. For operators on state spaces $V_i$ and $V_j$ (with $i \neq j$), the commutation relations are
\begin{align}
\hat{c}_i  \hat{c}_j = -  \hat{c}_j  \hat{c}_i \\ 
\hat{c}^\dagger_i  \hat{c}^\dagger_j = -  \hat{c}^\dagger_j  \hat{c}^\dagger_i \\ 
\hat{c}_i  \hat{c}^\dagger_j = -  \hat{c}^\dagger_j  \hat{c}_i  \label{eqn:ccdag_relation}
\end{align}
The first two of these are commonly summarized as:
\begin{align}
\{ \hat{c}_i, \hat{c}_j \} = 0 \nonumber \\
\{ \hat{c}^\dagger_i, \hat{c}^\dagger_j \} = 0
\end{align}
where $\{ a, b \} = a b + b a$ is the anti-commutator.

The third relation Eq.~(\ref{eqn:ccdag_relation}) has the special case that when $i=j$, $(\hat{c}_i \hat{c}^\dagger_i + \hat{c}^\dagger_i \hat{c}_i) = 1$. Therefore it is commonly summarized as 
\begin{align}
\{ \hat{c}_i, \hat{c}^\dagger_j \} = \delta_{ij} \ \ .
\end{align}

In many presentations, the anticommutation relations of the operators are stated as axioms, but they actually follow from the properties of the state space.
For example, 
\begin{align}
\hat{c}_1  \hat{c}_2   & =  \ket{0}_1 \otimes \bra{1}_1 \otimes \ket{0}_2 \otimes \bra{1}_2 \nonumber \\
& =  \ket{0}_1 \otimes (\bra{1}_1  \otimes \bra{1}_2) \otimes \ket{0}_2 \nonumber \\
& = \ket{0}_1 \otimes ( - \bra{1}_2  \otimes \bra{1}_1) \otimes \ket{0}_2 \nonumber \\
& = - \ket{0}_2 \otimes \bra{1}_2  \otimes \bra{0}_1 \otimes \ket{1}_1 \nonumber \\
& = - \hat{c}_2  \otimes \hat{c}_1 \ \ .
\end{align}

\section{Many-Body States and Gaussian States}

Basis states for many-body or many-particle systems can be defined or constructed as
\begin{align}
\ket{1_i 1_j 1_k} = \hat{c}^\dagger_i \hat{c}^\dagger_j \hat{c}^\dagger_k \ket{0} \ .
\end{align}
The most general many-particle state is a sum of these:
\begin{align}
\ket{\Psi} = \sum_{n_1, n_2, \ldots, n_N} \Psi^{n_1 n_2 \cdots n_N} \ket{n_1 n_2 \cdots n_N} \label{eqn:general_state}
\end{align}
where $\Psi^{n_1 n_2 \cdots n_N}$ is an arbitrary tensor of complex amplitudes. The anti-commuting nature of the fermions is automatically satisfied by the properties of the basis, and no additional symmetry property needs to be assumed for the amplitudes $\Psi^{n_1 n_2 \cdots n_N}$. 

A typical extra assumption, though, is that for a closed system the total number of particles $N_p$ is fixed. Mathematically this means $n_1 + n_2 + \ldots + n_N = N_p$ otherwise $\Psi^{n_1 n_2 \cdots n_N} = 0$. (Superconducting states are a common exception to this rule, though these are more properly understood in terms of density matrices or mixed states.)

A special subset of many-particle fermionic states are the \emph{Gaussian states}, more commonly known as \emph{Slater determinants}.
The general form of a Gaussian state with four fermions is:
\begin{align}
\ket{\Psi} & = \hat{A}^\dagger \hat{B}^\dagger \hat{C}^\dagger \hat{D}^\dagger \ket{0} \nonumber \\
& = \left( \sum_{i_1=1}^N A^{i_1} \hat{c}^\dagger_{i_1} \right) \left( \sum_{i_2=1}^N B^{i_2} \hat{c}^\dagger_{i_2} \right) \left( \sum_{i_3=1}^N C^{i_3} \hat{c}^\dagger_{i_3} \right) \left( \sum_{i_4=1}^N D^{i_4} \hat{c}^\dagger_{i_4} \right) \ket{0} \ \ \ .
\end{align}
Above we have introduced operators such as $\hat{A} =  \sum_{i_1=1}^N A^{i_1} \hat{c}^\dagger_{i_1}$. One describes the state above by saying that it contains four fermions which are in the ``orbitals'' $A, B, C$ and $D$. These orbitals are vectors $\vec{A}= [A^1 \ A^2 \ \ldots A^N]$ which are unit-normalized. The vectors are not required to be orthogonal, though the properties of the state are simpler to deduce if they are.

The Gaussian states are quite special and far from being a general fermionic state, even though they are elements of the general vector space of fermionic many-body states. To see that the Gaussian states are not general, let us write a two-particle Gaussian state 
$\ket{AB} = \hat{A}^\dagger \hat{B}^\dagger \ket{0}$ in the general form Eq.~(\ref{eqn:general_state}) for a fermionic state.
We see that for the state $\ket{AB}$, the amplitudes $\Psi^{n_1 n_2 \cdots n_N}$ take the form:
\begin{align}
\Psi^{0_1 0_2 \cdots 1_i \cdots 1_j \cdots 0_N} & = A_i B_j \ \  \text{ for}\ \sum_i n_i = 2\\
\Psi^{n_1 n_2 \cdots n_N} & = 0 \ \ \ \ \  \ \ \text{otherwise for}\ \sum_i n_i \neq 2 \ \ .
\end{align}
So the amplitude tensor $\Psi$ of $\ket{AB}$ is constructed from two vectors (it is rank-2 in a certain sense \footnote{The amplitude tensor for $\ket{AB}$ or a general two-particle Gaussian state is actually a bond-dimension 2 or rank 2 matrix product state}) and is not fully general.

To understand the origin of the name Slater determinant, consider the basis defined as $\ket{i} = \hat{c}^\dagger_i \ket{0}$, the so-called ``real-space basis''.
Expressing states in this basis is known as \emph{first quantization}.
We can express the amplitudes of $\ket{AB}$ in this basis, as $\psi_{AB}(i,j) = \braket{i,j|AB} = \bra{0} c^\dagger_i c^\dagger_j \ket{AB}$. One finds that
\begin{align}
\psi_{AB}(i,j) =  A^i B^j - B^i A^j  \label{eq:two_particle_slater}
\end{align}
such that the amplitudes in real space take the form of an antisymmetric function.
It was pointed out by Slater that expressions of the above type can be written as determinants of matrices whose columns are the orbital vectors $\vec{A}, \vec{B}, \ldots$ and whose rows are the spatial positions $i=1,2, \ldots, N$, hence the name Slater determinant. 

The first-quantized basis can be convenient for proving facts such as if $\vec{A} = \vec{B}$, then the state $\ket{AB} = 0$, as is evident from Eq.~(\ref{eq:two_particle_slater}) above. In the more general setting, the requirement that orbitals required be independent is known as the \emph{Pauli exclusion principle}. It is not a separate axiom about fermions, but follows from their vector space properties.


\section{Static Properties of Gaussian States}

\begin{itemize}
\item $[X]$ Define g operators.
\item $[\ \ \ ]$ Emphasize neither basis is preferred in this section. Just a linear mapping from one to the other.
\end{itemize}

In this section we will consider static (time-independent) properties of Gaussian states. To define a generic Gaussian state,
start by defining a matrix of orbitals $g^{j}\!_m$. This is assumed to be a unitary matrix $\bar{g}^{m'}\!_j g^j\!_m = \delta^{m'}_m$.
The $m^\text{th}$ column $[\psi^{(m)}]^{j} = g^{j}\!_m$ can be interpreted as the real-space wavefunction
for the single fermion that will be created by the $\hat{g}^\dagger_m$ operator defined next.

Define the operator $\hat{g}^\dagger_m$ and its conjugate as
\begin{align}
\hat{g}^\dagger_m & = \sum_j g^j\!_m \, \hat{c}^\dagger_j \\
\hat{g}^m & = \sum_j \bar{g}^m\!_j \, \hat{c}^j  \ .
\end{align}

The most generic form of a Gaussian state or Slater determinant is
\begin{align}
\ket{m_1 m_2 \ldots m_n}_g & = \hat{g}^\dagger_{m_1} \hat{g}^\dagger_{m_2} \cdots \hat{g}^\dagger_{m_n} \ket{0} \\
& = \left[ \prod_{m=1} (\hat{g}^\dagger_{m})^{\nu_m} \right] \ket{0}
\end{align}
where $\nu_m$ is the \emph{occupancy} of state $m$ in the $g$ basis. The occupancy is defined here as $\nu_m = 1$ if $m \in \{m_1, m_2, \ldots, m_n\}$ and $\nu_m = 0$ otherwise.

Now we will derive important results about properties of Gaussian states.

\subsection{Correlation matrix}

First consider the case of a Gaussian state 

Derivation strategy:
\begin{itemize}
\item Show that state like $\hat{c}^\dagger_{j_1} \hat{c}^\dagger_{j_2} |0\rangle$ has diagonal correlation matrix in c basis.
\item Argue for similar result in g basis, using canonical commutation relations.
\item Use canonical transformation to derive $\langle \mathbf{m} | \hat{c}^\dagger_i \hat{c}^j | \mathbf{m} \rangle$. 
\end{itemize}


\section{Dynamics of Gaussian States}

\section{Mathematical Results}


\textbf{Correlation functions and expectation values.}
\begin{itemize}

\item \emph{Correlation matrix:} Given a Gaussian state $\ket{g} = \hat{g}^\dagger_{m_1} \hat{g}^\dagger_{m_2} \cdots \hat{g}^\dagger_{m_n} \ket{0}$ where $\hat{g}^\dagger_m = g^{i}_m \hat{f}^\dagger_i$, the correlation matrix $C^j_i$ in the f basis is given by \textcolor{blue}{$C_i^j = \bra{g} \hat{f}^\dagger_i \hat{f}^j \ket{g} = \sum_n \bar{g}_i^n \nu_n g^j_n$}, where $\nu_m = \bra{g} \hat{n}_{m} \ket{g}$ is the occupancy of state $m$ in the $g$ basis ($\nu_m = 1$ if $m \in \{m_1, m_2, m_n\}$, $\nu_m = 0$ otherwise).

\item \emph{Same-basis correlation matrix:} Given a Gaussian state $\ket{f} = \hat{f}^\dagger_{j_1} \hat{f}^\dagger_{j_2} \cdots \hat{f}^\dagger_{j_n} \ket{0}$, the correlation matrix $C^j_i$ in the f basis is given by \textcolor{blue}{$C_i^j = \bra{f} \hat{f}^\dagger_i \hat{f}^j \ket{f} = \delta_j^i \nu_j$} where $\nu_j = \bra{f}\hat{n}_j\ket{f}$ is the occupancy of orbital $j$ in the $f$ basis.

\item \emph{Expectation value:} Given an operator $\hat{M} = m^{i}_j \hat{c}^\dagger_i c^j$ and a Gaussian state $\ket{f}$ with correlation matrix 
\mbox{$C_i^j = \bra{f} \hat{c}^\dagger_i \hat{c}^j \ket{f}$}, the expected value of $\hat{M}$ is \result{$\langle \hat{M} \rangle = \text{Tr}[mC]$}.

\end{itemize}


\textbf{Action of operators.}
\begin{itemize}

\item \emph{Action of creation operator:} Given state $\ket{g} = \hat{g}^\dagger_{m_1} \hat{g}^\dagger_{m_2} \cdots \hat{g}^\dagger_{m_n} \ket{0}$ where $\hat{g}^\dagger_m = g^{i}_m \hat{f}^\dagger_i$ and an operator $\hat{v}^\dagger = v^j \hat{f}^\dagger_i$, the correlation matrix of the state $\ket{g'} = \hat{v}^\dagger \ket{g}$ is given by \result{$C' = ||v_0||^2 C + v_0 v_0^\dagger$}, where $v_0 = (1-C)v$ is the projection of $v$ onto the unoccupied subspace of $C$.

\item \emph{Action of annihilation operator:} Given state $\ket{g} = \hat{g}^\dagger_{m_1} \hat{g}^\dagger_{m_2} \cdots \hat{g}^\dagger_{m_n} \ket{0}$ where $\hat{g}^\dagger_m = g^{i}_m \hat{f}^\dagger_i$ and an operator $\hat{w} = \bar{w}_i \hat{f}^i$, the correlation matrix of the state $\ket{g'} = \hat{w} \ket{g}$ is given by \result{$C' = ||w_1||^2 C - w_1 w_1^\dagger$}, where $w_1 = C w$ is the projection of $w$ onto the occupied subspace of $C$.


\end{itemize}


\textbf{Dynamics}. All of these results assume the Hamiltonian $\hat{H} = h^i_j \hat{c}^\dagger_i \hat{c}^j$ has a hopping matrix which can be diagonalized
as $h^i_j = \phi^i_n \epsilon_n \phi^n_j$. The ground state of $\hat{H}$ is $\ket{\Phi_0}$.
\begin{itemize}
\item \emph{Heisenberg time evolution of operators:} Define the operators $\hat{\phi}^\dagger_n = \sum_n \phi^j_n \hat{c}^\dagger_j$ and $\hat{\phi}^n = \sum_n \bar{\phi}_j^n \hat{c}^j$. These evolve in the Heisenberg picture as \result{$\hat{\phi}^n(t) = e^{i \hat{H} t} \hat{\phi}^n e^{-i \hat{H} t} = e^{- i\epsilon_n t} \hat{\phi}^n$} and \textcolor{blue}{$\hat{\phi}^\dagger_n(t) = e^{i \epsilon_n t} \hat{\phi}^\dagger_n$}.
\item \emph{Lesser Greens function:} Defined as $G^{<\, i}_j(t)= i \bra{\Phi_0} \hat{c}^\dagger_j \hat{c}^i(t) \ket{\Phi_0}$, it is given by \result{$G^{<\, i}_j(t) = i \sum_n \nu_n  \phi^i_n e^{-i \epsilon_n t}  \bar{\phi}_j^n $} where $\nu_n$ is the occupancy of the ground state in the $\phi$ basis.
\item \emph{Greater Greens function:} Defined as $G^{>\, i}_j(t) = -i \bra{\Phi_0} \hat{c}^i(t) \hat{c}^\dagger_j \ket{\Phi_0}$, it is given by \mbox{\result{$G^{>\, i}_j(t) = -i \sum_n (1-\nu_n) \phi^i_n e^{-i \epsilon_n t} \bar{\phi}_j^n $}} where $\nu_n$ is the occupancy of the ground state in the $\phi$ basis.
\item \emph{Retarded Greens function:} Defined as $G^{R\, i}_j(t) = \theta(t)(G^{>\, i}_j(t) - G^{<\, i}_j(t)) $, it is given by \result{$G^{R\, i}_j(t) = -i \theta(t) \sum_n \phi^i_n e^{-i \epsilon_n t} \bar{\phi}_j^n  = -i \theta(t) e^{-i h t}$}.
\item \emph{Propagator:} Define the propagator as $g^i_j(t) = i(G^{>\, i}_j(t) - G^{<\, i}_j(t))$. It can be written simply as \result{$g^i_j(t) = e^{-i h t} = \sum_n \phi^i_n e^{-i \epsilon_n t} \bar{\phi}_j^n$}.
\item \emph{Evolution of correlation matrix:} For a pure, Gaussian state $\ket{f}$ with correlation matrix $C_i^j = \bra{f} \hat{c}^\dagger_i \hat{c}^j \ket{f} = \bar{f}_i^n \nu_n f^j_n$, the orbitals evolve, in the Schr\"odinger picture, as $f^j_n(t) = g^i_j(t) f^j_n$. So \result{$C_i^j(t) =  \bar{g}^{i'}_i(t) C_{i'}^{j'}  g^j_{j'}(t)$}.

\end{itemize}


\end{document}
